\documentclass{article}

\usepackage{amsmath,graphicx,epsfig}
% \usepackage{epstopdf}
% \usepackage{euscript}
\usepackage{amsfonts}
\usepackage{amssymb}
\usepackage{float}

\usepackage[framemethod=tikz]{mdframed}

\usepackage{setspace}
% \doublespacing
\usepackage[skip=10pt plus1pt, indent=20pt]{parskip}
\emergencystretch=2em
% \parskip=1\baselineskip \advance\parskip by 0pt plus 2pt

% %%%%%%%%%%%%%%%%%%%%%%%%%%%%%%%%%%%%%%%%%%%%%%%%%%%%%%%%%%%%%%%%%%%%%%%%%%%%%%%%%%%%%%%%%%%%%%%%%%%%%%%%%%%%%%%%%
% We need hyperref for using LATEXDIFF
\usepackage{hyperref}

%DIF PREAMBLE EXTENSION ADDED BY LATEXDIFF
%DIF UNDERLINE PREAMBLE %DIF PREAMBLE
\RequirePackage[normalem]{ulem} %DIF PREAMBLE
\RequirePackage{color}\definecolor{RED}{rgb}{1,0,0}\definecolor{BLUE}{rgb}{0,0,1} %DIF PREAMBLE
\providecommand{\DIFaddtex}[1]{{\protect\color{blue}\uwave{#1}}} %DIF PREAMBLE
\providecommand{\DIFdeltex}[1]{{\protect\color{red}\sout{#1}}}                      %DIF PREAMBLE
%DIF SAFE PREAMBLE %DIF PREAMBLE
\providecommand{\DIFaddbegin}{} %DIF PREAMBLE
\providecommand{\DIFaddend}{} %DIF PREAMBLE
\providecommand{\DIFdelbegin}{} %DIF PREAMBLE
\providecommand{\DIFdelend}{} %DIF PREAMBLE
\providecommand{\DIFmodbegin}{} %DIF PREAMBLE
\providecommand{\DIFmodend}{} %DIF PREAMBLE
%DIF FLOATSAFE PREAMBLE %DIF PREAMBLE
\providecommand{\DIFaddFL}[1]{\DIFadd{#1}} %DIF PREAMBLE
\providecommand{\DIFdelFL}[1]{\DIFdel{#1}} %DIF PREAMBLE
\providecommand{\DIFaddbeginFL}{} %DIF PREAMBLE
\providecommand{\DIFaddendFL}{} %DIF PREAMBLE
\providecommand{\DIFdelbeginFL}{} %DIF PREAMBLE
\providecommand{\DIFdelendFL}{} %DIF PREAMBLE
%DIF HYPERREF PREAMBLE %DIF PREAMBLE
\providecommand{\DIFadd}[1]{\texorpdfstring{\DIFaddtex{#1}}{#1}} %DIF PREAMBLE
\providecommand{\DIFdel}[1]{\texorpdfstring{\DIFdeltex{#1}}{}} %DIF PREAMBLE
%DIF COLORLISTINGS PREAMBLE %DIF PREAMBLE
\RequirePackage{listings} %DIF PREAMBLE
\RequirePackage{color} %DIF PREAMBLE
\lstdefinelanguage{DIFcode}{ %DIF PREAMBLE
%DIF DIFCODE_UNDERLINE %DIF PREAMBLE
  moredelim=[il][\color{red}\sout]{\%DIF\ <\ }, %DIF PREAMBLE
  moredelim=[il][\color{blue}\uwave]{\%DIF\ >\ } %DIF PREAMBLE
} %DIF PREAMBLE
\lstdefinestyle{DIFverbatimstyle}{ %DIF PREAMBLE
	language=DIFcode, %DIF PREAMBLE
	basicstyle=\ttfamily, %DIF PREAMBLE
	columns=fullflexible, %DIF PREAMBLE
	keepspaces=true %DIF PREAMBLE
} %DIF PREAMBLE
\lstnewenvironment{DIFverbatim}{\lstset{style=DIFverbatimstyle}}{} %DIF PREAMBLE
\lstnewenvironment{DIFverbatim*}{\lstset{style=DIFverbatimstyle,showspaces=true}}{} %DIF PREAMBLE

\newcommand{\difDelContent}[1]{
  \DIFdelbegin \DIFdel{#1}\DIFdelend 
}
\newcommand{\difAddContent}[1]{
  \DIFaddbegin
  \begingroup\color{blue}#1\endgroup
  \DIFaddend
}

%DIF END PREAMBLE EXTENSION ADDED BY LATEXDIFF
% %%%%%%%%%%%%%%%%%%%%%%%%%%%%%%%%%%%%%%%%%%%%%%%%%%%%%%%%%%%%%%%%%%%%%%%%%%%%%%%%%%%%%%%%%%%%%%%%%%%%%%%%%%%%%%%%%

% %%%%%%%%%%%%%%%%%%%%%%%%%%%%%%%%%%%%%%%%%%%%%%%%%%%%%%%%%%%%%%%%%%%%%%%%%%%%%%%%%%%%%%%%%%%%%%%%%%%%%%%%%%%%%%%%%
% commands for simplifying the boxes
\newcommand{\refereeOverallComment}[2]{
  \noindent \textbf{Reviewer #1 gives their overall opinion as follows:}
  \begin{mdframed}[linecolor=black]
  #2
  \end{mdframed}
}


\newcommand{\refereeNoOneBox}[2]{
  \begin{mdframed}[linecolor=black]
  \noindent \textbf{Reviewer 1: (#1)} #2
  \end{mdframed}
}

\newcommand{\refereeNoTwoBox}[2]{
  \begin{mdframed}[linecolor=black]
  \noindent \textbf{Reviewer 2: (#1)} #2
  \end{mdframed}
}

\newcommand{\mdframedBox}[1]{
  \begin{mdframed}[linecolor=blue,topline=false,bottomline=false]
  \noindent #1
  \end{mdframed}
}

\newcommand{\response}[2]{
  \noindent \textbf{Response to (#1):} #2
}
% %%%%%%%%%%%%%%%%%%%%%%%%%%%%%%%%%%%%%%%%%%%%%%%%%%%%%%%%%%%%%%%%%%%%%%%%%%%%%%%%%%%%%%%%%%%%%%%%%%%%%%%%%%%%%%%%%

% %%%%%%%%%%%%%%%%%%%%%%%%%%%%%%%%%%%%%%%%%%%%%%%%%%%%%%%%%%%%%%%%%%%%%%%%%%%%%%%%%%%%%%%%%%%%
% SI unit
\usepackage{siunitx}
\DeclareSIUnit\torr{torr}
\DeclareSIUnit\ppm{ppm}
\DeclareSIUnit\ppb{ppb}
\DeclareSIUnit\bar{bar}
\DeclareSIUnit\gauss{G}
\DeclareSIUnit\neV{\nano\electronvolt}
\DeclareSIUnit\peV{\pico\electronvolt}
\DeclareSIUnit\kms{km/s}
\sisetup{print-unity-mantissa=false}  % To display only the exponent (i.e., 10⁻¹¹) 
\sisetup{separate-uncertainty = false} % use parenthesis style for uncertainty
% \sisetup{uncertainty-mode = separate} % use +/- style for uncertainty
\sisetup{uncertainty-mode = compact-marker} % turn this on to use style like 123.45(1.20)
% otherwise we  turn this on to use style like 123.45(120)
% %%%%%%%%%%%%%%%%%%%%%%%%%%%%%%%%%%%%%%%%%%%%%%%%%%%%%%%%%%%%%%%%%%%%%%%%%%%%%%%%%%%%%%%%%%%%

\usepackage{changes}

% %%%%%%%%%%%%%%%%%%%%%%%%%%%%%%%%%%%%%%%%%%%%%%%%%%%%%%%%%%%%%%%%%%%%%%%%%%%%%%%%%%%%%%%%%%%%
% for chemical forumals
\usepackage{chemformula}
% %%%%%%%%%%%%%%%%%%%%%%%%%%%%%%%%%%%%%%%%%%%%%%%%%%%%%%%%%%%%%%%%%%%%%%%%%%%%%%%%%%%%%%%%%%%%


\begin{document}

\title{Response to Reviewer 1 for Manuscript ID universe-4542087}
%==================================

\date{\today}
\author{Yuzhe Zhang}

\maketitle

\noindent Dear Reviewer,
 
I sincerely appreciate your careful evaluation.
Your comments have identified important points and provided valuable suggestions.
Below, I provide responses to your comments and describe the related changes to the manuscript.

\refereeOverallComment{1}{
    The article  presents an interesting tool and is written in a very clear manner. However the article before publishing needs a major revision, particulary in the calibration section where findings should be clearly presented and possibly compared with real measurements.

}

I am very glad to see the positive evaluation and constructive comments from the referee. 
I fully respect your opinion and have taken action to improve the manuscript according to your suggestions. 
Edits of the manuscript are listed below. 

\section{Enriched examples of calibrations}
The examples of calibrations have been enriched in the revised manuscript in the following ways. 


\section{Detailed description of the quantitative calibration}

\section{Comparison with real measurements}

In the calibration section, the simulation results are compared with the well tested theoretical predictions. 
What can even be more convincing is to compare the simulation results with real measurements, as suggested by the referee.
Therefore, I have added examples of real measurements in the appendix to demonstrate the validity of the simulation results.

In Appendix A, I have added two examples of real measurements. 
The first example is the spin-echo NMR simulation using a CPMG pulse train, which is compared with the real measurements in Ref. 
\mdframedBox{}

As can be seen, the simulation results are in good agreement with the real measurements, which validates the accuracy of the simulation tool. 

\section{Further description of multi-thread implementation}
% \mdframedBox{
% \difDelContent{It opens the door to probing large swaths of hitherto unexplored mass-coupling parameter space in the future by using hyperpolarized samples.}

% \difAddContent{
% CASPEr-Gradient Low-Field will probe the mass range from 
% $\SI{4.1}{peV/c^2}$ to $\SI{17}{neV/c^2}$
% with hyperpolarized samples to boost the sensitivity beyond the astronomical limits.
% }
% }

~\
~\

I hope that the responses adequately address your comments and that the revisions are sufficient.
I believe that the manuscript has been substantially improved and hope that you will find it suitable for publication in its present form.
If any errors remain or if any of the comments have not been fully addressed, please let me know.
I would be pleased to make further improvements.

\noindent Yours sincerely,


\noindent Yuzhe Zhang

\bibliographystyle{unsrt}
\bibliography{../references}

\end{document} 
